% MatID: calculo_1
% BookID: cd1_b4
\documentclass[12pt]{book}
\usepackage[utf8]{inputenc}
\usepackage{amsmath, amssymb, amsthm}
\usepackage{graphicx}
\usepackage{enumitem}

% Teoremas y entornos
\newtheorem{theorem}{Teorema}[section]
\newtheorem{lemma}[theorem]{Lema}
\theoremstyle{definition}
\newtheorem{definition}{Definición}[section]
\newtheorem{example}{Ejemplo}[section]

\title{Introducción al Cálculo y al Análisis Matemático}
\author{Richard Courant y Fritz John}
\date{Volumen 5}

\begin{document}

\maketitle
\tableofcontents

\chapter{Introducción}

Desde la antigüedad las nociones intuitivas de cambio continuo, crecimiento y movimiento, han atraído la atención de las mentes científicas. Empero, el camino hacia el entendimiento de variación continua no fue logrado sino hasta el siglo XVII, cuando la ciencia moderna surgió y se desarrolló rápidamente en cercana conjunción con el cálculo integral y diferencial, llamado brevemente Cálculo, y el análisis matemático.

Las nociones básicas del Cálculo son las de derivada e integral: la derivada es una medida de la ``rapidez'' de una variación; la integral, una medida del efecto total de un proceso de cambio continuo. Una comprensión precisa de estos conceptos y de su abrumadora fecundidad descansa sobre los conceptos de límite y función; los cuales a su vez dependen de un entendimiento del continuo de números. Sólo gradualmente, penetrando más y más en la esencia del Cálculo, puede uno apreciar su potencia y belleza. En este capítulo introductorio se explicarán los conceptos básicos de número, función y límite, primero de manera simple e intuitiva, y después mediante argumentos cuidadosos.

\section{El continuo de números}

Los enteros positivos o números naturales $1, 2, 3, \dots$ son símbolos abstractos para indicar ``cuántos'' objetos hay en una colección o conjunto de elementos discretos.

Estos símbolos están despojados de toda referencia a las cualidades concretas de los objetos contados, ya sea que se trate de personas, átomos, casas o cualesquiera objetos.

Los números naturales son el instrumento adecuado para contar elementos de una colección o ``conjunto''. Sin embargo, no bastan para otro objetivo igualmente importante: medir cantidades tales como la longitud de una curva y el volumen o peso de un cuerpo. La pregunta ``¿cuánto?'' no puede ser contestada inmediatamente en términos de los números naturales. La profunda necesidad de expresar medidas de cantidades en términos de lo que nosotros quisiéramos llamar números nos obliga a extender el concepto de número de manera que podamos describir una graduación continua de las medidas. Esta extensión es llamada el continuo de números o el sistema de ``números reales'' (un nombre no descriptivo pero aceptado generalmente). La extensión del concepto de número al de continuo es de tal manera convincente que fue utilizada por todos los grandes matemáticos y científicos de épocas anteriores sin hacer preguntas escudriñadoras. No fue sino hasta el siglo XIX que los matemáticos se vieron impulsados a buscar un fundamento lógico más firme para el sistema de los números reales. La subsiguiente formulación precisa de estos conceptos condujo, a su vez, a un progreso mayor en matemáticas. Principiaremos con un tratamiento intuitivo sin complicaciones, y posteriormente se dará un análisis más profundo del sistema de los números reales.

\subsection{El sistema de números naturales y su extensión. Numeración y medición}

Los números naturales y los racionales. La sucesión de números ``naturales'' $1, 2, 3, \dots$ se considera como dada. No se requiere discutir cómo estos entes abstractos, los números, pueden ser catalogados desde un punto de vista filosófico. Para el matemático, y para cualquiera que trabaje con números, es importante conocer solamente las reglas o leyes mediante las cuales éstos pueden combinarse para obtener otros números naturales. Estas leyes forman la base de las reglas conocidas para sumar y multiplicar números en el sistema decimal; incluyen las leyes conmutativas
\[
a + b = b + a \quad \text{y} \quad ab = ba,
\]
las leyes asociativas
\[
a + (b + c) = (a + b) + c \quad \text{y} \quad a(bc) = (ab)c,
\]
la ley distributiva
\[
a(b + c) = ab + ac,
\]
la ley de la cancelación, esto es que $a + c = b + c$ implica $a = b$, etc.

Las operaciones inversas, substracción y división, no son siempre posibles dentro del conjunto de los números naturales; no podemos restar $2$ de $1$ o dividir $1$ entre $2$ y permanecer dentro de ese conjunto. Para hacer posibles estas operaciones sin restricción, nos vemos obligados a extender el concepto de número inventando el número $0$, los enteros ``negativos'' y las fracciones. La totalidad de todos estos números se llama la clase o conjunto de números racionales; todos ellos se obtienen de la unidad utilizando las ``operaciones racionales'' de cálculo, a saber: la adición, substracción, multiplicación y división.

Un número racional puede expresarse siempre en la forma $p/q$, donde $p$ y $q$ son enteros y $q \neq 0$. Podemos hacer que esta representación sea única exigiendo que $q$ sea positivo y que $p$ y $q$ no tengan un factor común mayor que $1$.

Dentro del dominio de los números racionales todas las operaciones racionales, adición, multiplicación, substracción y división (excepto la división por cero), pueden efectuarse y producir nuevamente números racionales. Como se sabe de la aritmética elemental, las operaciones con números racionales obedecen las mismas leyes que las operaciones con números naturales; por ello los números racionales extienden el sistema de enteros positivos en una forma completamente natural.

\textbf{Representación gráfica de los números racionales.} Los números racionales son usualmente representados en forma gráfica mediante puntos en una línea recta $L$, el eje numérico (o recta numérica). Tomando un punto arbitrario de $L$ como el origen o punto $0$ y otro punto arbitrario como el punto $1$, utilizamos la distancia entre estos dos puntos como escala o unidad de medida y definimos la dirección de $0$ a $1$ como ``positiva''. La línea recta con una dirección así impuesta se llama recta dirigida. Es costumbre dibujar a $L$ de manera que el punto $1$ se encuentre a la derecha del punto $0$ (Figura 1.1). La localización de cualquier punto $P$ en $L$ está completamente determinada por dos partes de información: la distancia de $P$ desde el origen $0$ y la dirección de $0$ a $P$ (a la derecha o izquierda del $0$). El punto $P$ en $L$ que representa un número racional positivo $x$ se encuentra a una distancia de $x$ unidades a la derecha del $0$. Un número racional negativo $x$ es representado por el punto $-x$ unidades a la izquierda del $0$. En cualquier caso la distancia del $0$ al punto que representa a $x$ se llama el valor absoluto de $x$, denotado por $|x|$, y se tiene
\[
|x| = \begin{cases}
x & \text{para } x \ge 0 \\
-x & \text{para } x < 0
\end{cases}
\]
Se observa que $|x|$ nunca es negativo y es igual a cero sólo cuando $x = 0$.

Recordemos de la geometría elemental que con regla y compás es posible subdividir la unidad de longitud en un número cualquiera de partes iguales. Se sigue que cualquier longitud racional puede ser construida, y, por lo tanto, que el punto que representa un número racional $x$ puede ser encontrado por métodos puramente geométricos.

De esta forma se obtiene una representación geométrica de números racionales mediante puntos en $L$, los llamados puntos racionales. Conforme a nuestra notación para los puntos $0$ y $1$, nos tomamos la libertad de denotar tanto los números racionales como los puntos correspondientes en $L$ por el mismo símbolo $x$.

La relación $x < y$ para dos números racionales significa geométricamente que el punto $x$ se encuentra a la izquierda del punto $y$. En ese caso la distancia entre los puntos es $y - x$ unidades. Si $x > y$, la distancia es $x - y$ unidades. En cualquier caso la distancia entre dos puntos racionales $x, y$ de $L$ es $|y - x|$ unidades y es también un número racional.

\begin{figure}[h]
\centering
% \includegraphics{figura1.2.png}
\caption{Subdivisión del eje numérico.}
\end{figure}

Un segmento en $L$ con puntos extremos $a, b$ con $a < b$ se denominará intervalo. El segmento particular con puntos extremos $0, 1$ se llama intervalo unitario. Si los puntos extremos están incluidos en el intervalo, se dirá que el intervalo es cerrado; si los puntos extremos se excluyen, el intervalo se llama abierto. El intervalo abierto, denotado por $(a, b)$, consiste en aquellos puntos $x$ para los cuales $a < x < b$, esto es, en aquellos puntos que se encuentran ``entre'' $a$ y $b$. El intervalo cerrado, denotado por $[a, b]$, consiste en los puntos $x$ para los cuales $a \le x \le b$. En cualquier caso la longitud del intervalo es $b - a$.

Los puntos que corresponden a los enteros $0, \pm 1, \pm 2, \dots$ subdividen el eje numérico en intervalos de longitud unitaria. Todo punto en $L$ es un punto extremo o bien un punto interior de uno de los intervalos de la subdivisión. Si se sigue subdividiendo todo intervalo en $q$ partes iguales, se obtiene una subdivisión de $L$ en intervalos de longitud $1/q$ por puntos racionales de la forma $p/q$. Todo punto $P$ de $L$ es entonces un punto racional de la forma $p/q$, o bien se encuentra entre dos puntos racionales sucesivos $p/q$ y $(p+1)/q$. Puesto que puntos sucesivos de la subdivisión distan $1/q$ unidades, se sigue que podemos encontrar un punto racional $p/q$ cuya distancia desde $P$ no excede $1/q$ unidades. El número $1/q$ puede hacerse tan pequeño como se desee escogiendo $q$ como un entero positivo suficientemente grande. Por ejemplo, escogiendo $q = 10^n$ (donde $n$ es cualquier número natural), podemos encontrar una ``fracción decimal'' $x = p/10^n$ cuya distancia desde $P$ es menor que $1/10^n$. Aunque no se afirma que todo punto de $L$ es un punto racional, se ve por lo menos que pueden encontrarse puntos racionales arbitrariamente próximos a cualquier punto $P$ de $L$.

\textbf{Densidad}. La proximidad arbitraria de puntos racionales a un punto $P$ dado de $L$ se expresa diciendo: Los puntos racionales son densos en el eje numérico. Es claro que aun conjuntos menores de números racionales son densos, por ejemplo, los puntos $x = p/10^n$, para todos los números naturales $n$ y enteros $p$.

La densidad implica que entre dos puntos racionales cualesquiera $a$ y $b$ existe una infinidad de otros puntos racionales. En particular, el punto medio entre $a$ y $b$, $c = \frac{1}{2}(a + b)$, correspondiente a la media aritmética de los números $a$ y $b$, es nuevamente racional. Tomando los puntos medios de $a$ y $c$, de $b$ y $c$, y continuando de esta manera, se puede obtener cualquier número de puntos racionales entre $a$ y $b$.

Un punto arbitrario $P$ en $L$ puede ser localizado con cualquier grado de precisión utilizando puntos racionales. Puede parecer entonces a primera vista que la labor de localizar $P$ por un número ha sido lograda introduciendo los números racionales. Después de todo, en la realidad física las cantidades nunca se dan o conocen con absoluta precisión sino que siempre con sólo un grado de incertidumbre y por ello pueden muy bien ser consideradas como medidas por números racionales.

\textbf{Cantidades inconmensurables}. Los números racionales, aunque densos, no bastan como base teórica de medición por medio de números. Dos cantidades cuya razón es un número racional se llaman conmensurables, debido a que pueden ser expresadas como múltiplos enteros de una unidad común. Desde época tan remota como la de los siglos V o VI a.C. los matemáticos y filósofos griegos hicieron el sorprendente y profundo descubrimiento de que existen cantidades que no son conmensurables con una unidad dada. En particular, existen segmentos de recta que no son múltiplos racionales de un segmento unitario dado.

Es fácil dar un ejemplo de una longitud inconmensurable con la longitud unidad: la diagonal $l$ de un cuadrado con lados de longitud unidad. Pues, por el teorema de Pitágoras, el cuadrado de esta longitud $l$ debe ser igual a $2$. Por ello, si $l$ fuese un número racional y consecuentemente igual a $p/q$, donde $p$ y $q$ son enteros positivos, tendríamos $p^2 = 2q^2$. Podemos suponer que $p$ y $q$ no tienen factores comunes, pues tales factores comunes podrían ser cancelados en primer lugar. De acuerdo con la ecuación anterior, $p^2$ es un número par; de aquí que $p$ mismo debe ser par, digamos $p = 2p'$. Sustituyendo $p$ por $2p'$ se obtiene $4p'^2 = 2q^2$, o bien $q^2 = 2p'^2$; consecuentemente, $q^2$ sería par y por lo tanto $q$ también par. Esto es lo mismo que decir, que $p$ y $q$ tienen ambos el factor $2$. Sin embargo, esto contradice nuestra hipótesis de que $p$ y $q$ no tienen un factor común. Puesto que la suposición de que la diagonal puede ser representada por una fracción $p/q$ da lugar a una contradicción, la suposición es falsa. Este razonamiento, ejemplo característico de prueba indirecta, muestra que el símbolo $\sqrt{2}$ no puede corresponder a ningún número racional. Otro ejemplo es $\pi$, la razón de la circunferencia de un círculo a su diámetro. La demostración de que $\pi$ no es racional es mucho más complicada y se obtuvo solamente en tiempos modernos (Lambert, 1761). Es fácil encontrar muchas cantidades inconmensurables; de hecho, las cantidades inconmensurables son, en cierto sentido, mucho más comunes que las conmensurables.

\textbf{Números irracionales}. Debido a que el sistema de números racionales no es suficiente para la geometría, es necesario inventar nuevos números como medidas de cantidades inconmensurables: estos nuevos números se llaman ``irracionales''. Los antiguos griegos no pusieron de relieve el concepto de número abstracto; sin embargo, consideraron entes geométricos, tales como segmentos de recta, como los elementos básicos. De manera puramente geométrica ellos desarrollaron un sistema lógico para tratar y operar con cantidades inconmensurables, así como también con las conmensurables (racionales). Este logro tan importante, iniciado por los pitagóricos, fue desarrollado grandemente por Eudoxio y está expresado extensamente en el famoso \textit{Elementos} de Euclides. En los tiempos modernos las matemáticas fueron creadas de nuevo y vastamente desarrolladas sobre el fundamento de conceptos numéricos en lugar de conceptos geométricos. Con la introducción de la geometría analítica, se desarrolló una reversión de orientación e importancia en la relación antigua entre números y cantidades geométricas, y así la teoría clásica de los inconmensurables fue poco menos que olvidada u omitida. Se supuso como un hecho el que a todo punto del eje numérico le corresponde un número racional o irracional y que esta totalidad de números ``reales'' obedece las mismas leyes aritméticas que los números racionales. Sólo posteriormente, en el siglo XIX, se sintió la necesidad de justificar tal hipótesis y finalmente fue satisfecha completamente en un notable folleto de Dedekind cuya lectura es fascinante aún hoy en día.\footnote{R. Dedekind, ``Nature and Meaning of Number'', en \textit{Essays on Number}, Londres y Chicago, 1901. (El primero de estos ensayos, ``Continuity and Irrational Numbers'', proporciona una información detallada de la definición y leyes de operación con números reales.) Reimpreso bajo el título \textit{Essays on the Theory of Numbers}, Dover, Nueva York, 1964.}

En efecto, Dedekind mostró que el tratamiento intuitivo e ``ingenuo'' seguido por todos los grandes matemáticos desde Fermat y Newton hasta Gauss y Riemann estaba sobre la pista correcta: que el sistema de números reales (como símbolos para las longitudes de segmentos, o definidos de alguna otra manera) es un instrumento coherente y completo para la medición científica, y que en este sistema las reglas de cálculo del sistema de números racionales siguen siendo válidas.

Sin perjuicio, podríamos dejarlo así y dirigirnos directamente a la esencia del cálculo. Sin embargo, para un entendimiento más profundo del concepto de número real, el cual es necesario para nuestro trabajo posterior, la siguiente explicación, así como el Suplemento a este capítulo, deberían ser estudiados.

\subsection{Números reales e intervalos encajados}

Por el momento imaginemos los puntos en una recta $L$ como los elementos básicos del continuo. Postulamos el que a cada punto en $L$ le corresponde un ``número real'' $x$, su coordenada, y que para estos números $x$, $y$ las relaciones recién descritas para los números racionales conservan su significado. En particular, la relación $x < y$ indica orden en $L$ y la expresión $|y - x|$ significa la distancia entre el punto $x$ y el punto $y$. El problema básico es el de relacionar estos números (o mediciones en el continuo de puntos dado geométricamente) con los números racionales considerados originalmente y, por tanto, finalmente con los enteros. Además, debemos explicar cómo operar con los elementos de este ``continuo de números'' en la misma forma que con los números racionales. Con el tiempo, formularemos el concepto del continuo independientemente de los conceptos geométricos intuitivos, pero por ahora postergamos algo de la discusión más abstracta hasta el Suplemento.

¿Cómo podemos describir un número real irracional? Para algunos números, tales como $\sqrt{2}$ o $\pi$, podemos dar una caracterización geométrica simple, pero esto no es siempre factible. Un método suficientemente flexible para obtener cualquier punto real $x$ consiste en describir el valor $x$ por una sucesión de aproximaciones racionales de precisión cada vez mayor. Específicamente, aproximaremos $x$ simultáneamente desde la derecha y desde la izquierda incrementando sucesivamente la precisión y de tal manera que el margen de error se aproxime a cero. En otras palabras, utilizamos una ``sucesión'' de intervalos racionales que contienen a $x$, con cada intervalo de la sucesión conteniendo al siguiente de tal manera que la longitud del intervalo, y con él el error de aproximación, puede hacerse, tomando intervalos suficientemente alejados en la sucesión, menor que cualquier número positivo especificado. Para comenzar, confinese $x$ a un intervalo cerrado $I_1 = [a_1, b_1]$, esto es, $a_1 \le x \le b_1$ donde $a_1$ y $b_1$ son racionales. Dentro de $I_1$ consideramos un ``subintervalo'' $I_2 = [a_2, b_2]$ que contiene a $x$, esto es, $a_1 \le a_2 \le x \le b_2 \le b_1$ donde $a_2$ y $b_2$ son racionales. Por ejemplo, podemos escoger para $I_2$ una de las mitades de $I_1$, pues $x$ debe encontrarse en uno o ambos de los semiintervalos. Dentro de $I_2$ consideramos un subintervalo $I_3 = [a_3, b_3]$ que también contiene a $x$: $a_1 \le a_2 \le a_3 \le x \le b_3 \le b_2 \le b_1$ donde $a_3$ y $b_3$ son racionales, etc. Se requiere que la longitud del intervalo $I_n$ tienda a cero con $n$ creciente; esto es, que la longitud de $I_n$ sea menor que cualquier número positivo preasignado, para toda $n$ suficientemente grande. Un conjunto de intervalos cerrados $I_1, I_2, I_3, \dots$, cada uno conteniendo al siguiente y tal que las longitudes tiendan a cero se denominará ``sucesión anidada de intervalos'' o ``encaje'' de intervalos. El punto $x$ está determinado de manera única por la sucesión o encaje; esto es, ningún otro punto $y$ puede pertenecer a todo $I_n$, puesto que la distancia entre $x$ e $y$ excede la longitud de $I_n$ una vez que $n$ es suficientemente grande. Puesto que aquí escogemos siempre puntos racionales para los puntos extremos de los $I_n$ y puesto que todo intervalo con puntos extremos racionales es descrito por dos números racionales, vemos que todo punto $x$ de $L$, esto es, todo número real, puede ser precisamente determinado con la ayuda de una infinidad de números racionales. La aseveración inversa no es tan obvia; la aceptaremos como un axioma básico.

\begin{definition}[Postulado de los encajes de intervalos]
Si $I_1, I_2, I_3, \dots$ forman una sucesión ``encaje'' de intervalos con puntos extremos racionales, existe un punto $x$ contenido en todo $I_n$.
\end{definition}

Es importante destacar que para un encaje los intervalos $I_n$ son cerrados. Si, por ejemplo, $I_n$ denota el intervalo abierto $0 < x < 1/n$, entonces cada $I_n$ contiene al siguiente y las longitudes de los intervalos tienden a cero, pero no existe $x$ contenido en todo $I_n$.

\subsection{Fracciones decimales. Bases distintas a la de diez}

La representación decimal de los números reales es un caso especial del método de intervalos encajados. En el sistema decimal se subdivide el intervalo unitario en $10$ partes iguales y se continúa subdividiendo cada subintervalo en $10$ partes iguales, etc. Si se desea describir un número real $x$ por este método, se determina primero el entero $c_0$, escogiendo $c_0$ como el mayor entero que no excede a $x$ (esto es, $c_0 = [x]$ es la parte entera de $x$). Entonces $x$ pertenece al intervalo $I_0 = [c_0, c_0 + 1]$. Este intervalo se divide en $10$ partes iguales, determinando los puntos $c_0 + 1/10, c_0 + 2/10, \dots, c_0 + 9/10$. El número $x$ debe pertenecer a uno, y por lo menos, de estos subintervalos. Más aún, si $x$ es precisamente uno de los puntos de subdivisión (por ejemplo, $x = c_0 + 3/10$), entonces pertenece a dos subintervalos adyacentes. Para hacer la determinación única, se conviene en tomar el subintervalo cerrado a la derecha, digamos $[c_0 + 3/10, c_0 + 4/10]$. De esta forma se asigna un entero $c_1$ (entre $0$ y $9$) tal que
\[
c_0 + \frac{c_1}{10} \le x \le c_0 + \frac{c_1+1}{10}.
\]
Entonces se subdivide este intervalo en $10$ partes iguales, determinando un dígito $c_2$ ($0 \le c_2 \le 9$) tal que
\[
c_0 + \frac{c_1}{10} + \frac{c_2}{10^2} \le x \le c_0 + \frac{c_1}{10} + \frac{c_2+1}{10^2},
\]
y así sucesivamente. De esta manera se obtiene una sucesión de dígitos $c_1, c_2, c_3, \dots$ correspondientes al número real $x$. En este proceso todos excepto el primero son dígitos que tienen valores desde cero hasta nueve solamente. En la notación decimal ordinaria, la conexión entre $x$ y $c_0, c_1, c_2, \dots$ se indica escribiendo
\[
x = c_0 + 0.c_1c_2c_3\ldots
\]
(Usualmente el entero $c_0$ mismo es también escrito en notación decimal si $c_0$ es positivo.) Inversamente, por el axioma de continuidad, toda aquella expresión que denota una fracción decimal infinita representa un número real. Es posible que existan dos representaciones decimales diferentes del mismo número; por ejemplo, $1 = 0.99999\ldots = 1.00000\ldots$ En nuestra construcción el entero $c_0$ está determinado de manera única por $x$, a menos que el propio $x$ sea un entero. En este caso podríamos elegir $c_0 = x$ o bien $c_0 = x-1$. Una vez efectuada la elección, el dígito $c_1$ es único a menos que $x$ sea uno de los nuevos puntos que subdividen $I_0$ en diez partes iguales. Continuando así, encontramos que $c_0$ y todos los $c_k$ están determinados de manera única por $x$ a menos que $x$ resulte un punto de subdivisión en algún paso. Si esto sucediera por primera vez en el $n$-ésimo paso, entonces
\[
x = c_0 + \frac{c_1}{10} + \cdots + \frac{c_n}{10^n}
\]
donde $c_1, c_2, \dots, c_n$ son dígitos y donde $c_n > 0$, puesto que de otra manera $x$ hubiera sido un punto de subdivisión en algún paso anterior. Se sigue que $I_{n+1}$ es el intervalo $[x, x + 1/10^{n+1}]$ o bien el intervalo $[x - 1/10^{n+1}, x]$. En el primer caso $x$ será el punto extremo izquierdo de todos los intervalos posteriores $I_{n+2}, I_{n+3}, \dots$, y en el segundo caso, el punto extremo derecho. Esto nos lleva entonces a la representación decimal
\[
x = c_0 + 0.c_1c_2\ldots c_n000\ldots
\]
o bien a la representación
\[
x = c_0 + 0.c_1c_2\ldots (c_n-1)99999\ldots
\]
De aquí que el único caso en el cual puede surgir una ambigüedad es para un número racional $x$ que pueda ser escrito como una fracción que tiene una potencia de diez para su denominador. Podemos aún eliminar esta ambigüedad excluyendo representaciones decimales en las cuales todos los dígitos a partir de cierto punto son nueves.

En la representación decimal de números reales el papel especial desempeñado por el número diez es puramente incidental. La única razón evidente para el extenso uso del sistema decimal es la facilidad de contar por decenas con nuestros dedos (dígitos). Todo entero $p$ mayor que uno puede servir igualmente bien. Podrían utilizarse $p$ subdivisiones iguales en cada paso. Un número real $x$ sería entonces representado en la forma
\[
x = c_0 + 0.c_1c_2c_3\cdots,
\]
donde $c_0$ es un entero, y ahora $c_1, c_2, \dots$ tienen uno de los valores $0, 1, 2, \dots, p-1$. Esta representación caracteriza nuevamente a $x$ por un encaje de intervalos, a saber
\[
c_0 + \frac{1}{p}c_1 + \cdots + \frac{1}{p^n}c_n \le x \le c_0 + \frac{1}{p}c_1 + \cdots + \frac{1}{p^n}c_n + \frac{1}{p^n}.
\]
Si $x$ es positivo o cero, el entero $c_0$ es también positivo o cero, y el propio $c_0$ tiene un desarrollo finito de la forma
\[
c_0 = d_0 + pd_1 + p^2d_2 + \cdots + p^kd_k,
\]
donde $d_0, d_1, \dots, d_k$ toman uno de los valores $0, 1, \dots, p-1$. La representación completa de $x$ ``en la base $p$'' toma la forma
\[
x = d_k d_{k-1} \cdots d_1 d_0 . c_1 c_2 c_3 \cdots.
\]
Si $x$ es negativo podemos utilizar este tipo de representación para $-x$.

\begin{figure}[h]
\centering
% \includegraphics{figura1.4.png}
\caption{La fracción $2\frac{3}{4}$ en el sistema binario.}
\end{figure}

Bases distintas a la de $10$ han sido utilizadas extensamente en la actualidad. Siguiendo la dirección de los antiguos babilonios, los astrónomos durante muchos siglos representaron sistemáticamente los números como fracciones ``sexagesimales'', con $p = 60$ como la base.

\textbf{Representación binaria.} El sistema ``binario'' con la base $p = 2$ tiene un especial interés teórico y es útil en el diseño lógico de máquinas calculadoras electrónicas. En el sistema binario los dígitos tienen solamente dos valores posibles, el cero y el uno. El número $2\frac{3}{4}$, por ejemplo, sería escrito $10.11$, de acuerdo con la fórmula
\[
\frac{11}{4} = 2^2 \cdot 1 + 2^1 \cdot 0 + 1 \cdot 1 + \frac{1}{2} \cdot 1 + \frac{1}{2^2} \cdot 1.
\]

\textbf{Cálculo con números reales.} Aunque la definición de números reales y sus representaciones decimal infinita, binaria, etc., son naturales, puede no parecer obvio el que uno pueda operar con el continuo de números exactamente como con números racionales, efectuando las operaciones racionales y reteniendo las leyes de la aritmética, tales como las leyes asociativa, conmutativa y distributiva. La demostración es simple aunque algo tediosa. En lugar de dificultar el camino hacia la esencia viva del análisis considerando la cuestión aquí, aceptaremos temporalmente la posibilidad de cálculo aritmético ordinario con los números reales. Una comprensión más profunda de la estructura lógica subyacente al concepto de número será alcanzada cuando se descubra la idea de límite y sus implicaciones. (Ver el Suplemento a este capítulo, p. 112.)

\subsection{Definición de vecindad o entorno}

No sólo las operaciones racionales, sino también relaciones de orden o desigualdades para números reales obedecen las mismas reglas que los números racionales.

Pares de números reales $a$ y $b$ con $a < b$ dan nuevamente lugar a intervalos cerrados $[a, b]$ (dados por $a \le x \le b$) e intervalos abiertos $(a, b)$ (dados por $a < x < b$). Frecuentemente tendremos necesidad de asociar con un punto $x_0$ los diferentes intervalos abiertos que contienen a ese punto o bien, específicamente, tenerlo como centro; a esos intervalos los llamaremos vecindades o entornos del punto. Más precisamente, para todo positivo $\varepsilon$, la $\varepsilon$-vecindad del punto $x_0$ consiste de los valores $x$ para los cuales $x_0 - \varepsilon < x < x_0 + \varepsilon$; esto es, es el intervalo $(x_0 - \varepsilon, x_0 + \varepsilon)$. Todo intervalo abierto $(a, b)$ que contiene a un punto $x_0$, también contiene a una vecindad completa de $x_0$.

Habiendo definido intervalos con puntos extremos reales, podemos formar ahora encajes de intervalos utilizando la misma definición como en el caso de puntos extremos racionales. Es de lo más importante para la coherencia lógica del cálculo que para todo encaje de intervalos con puntos extremos reales exista un número real contenido en todos ellos. (Ver Suplemento, p. 118.)

\subsection{Desigualdades}

\subsubsection{Reglas básicas}

Las desigualdades juegan un papel mucho mayor en matemáticas superiores que en matemáticas elementales. A menudo el valor preciso de una cantidad $x$ es difícil de determinar, mientras que puede ser fácil efectuar una estimación de $x$, esto es, mostrar que $x$ es mayor que alguna cantidad conocida $a$ y menor que alguna otra cantidad $b$. Para muchos propósitos sólo la información contenida en tal estimación de $x$ es significativa. Recordaremos por ello brevemente algunas de las reglas elementales acerca de las desigualdades.

El
\end{document}
